\chapter{Introduction and Background}

We first fix some notations and background. Let \(k\) be an algebraically closed base field (\(k=\bar{k}\)). All rings are \(k\)-algebras, all schemes are \(k\)-schemes and all additive categories are \(k\)-linear. Let \(A\) be a ring and denote the category of right \(A\)-modules by \(\Mod(A)\). Similarly, if \(R=\oplus_{i\in \mathbb{Z}}R_i\) is a \(\mathbb{Z}\)-graded ring.

\begin{definition}[Grothendieck category]
     Let \(\mathcal{A}\) be an abelian category. We say \(\mathcal{A}\) is a Grothendieck category if the following is satisfied:
     \begin{itemize}
      \item \(\mathcal{A}\) has all arbitrary direct sums.
      \item Direct limits of short exact sequence are exact. 
      \item \mathcal{A} has a generator \(G\), that is \(\hom_{\mathcal{A}}(G,-)\) is a faithful functor. Or equivalently if all coproducts exists, for all object \(X\) in \(\mathcal{A}\), we have a epimorphism 
      \[\bigoplus_{i \in I}G\twoheadrightarrow X.\]
     \end{itemize}
\end{definition}

We first reformulate the well-known results in the commutative case. First we consider the affine case. Let \(X=\spec A\) be an affine scheme. A classical result tells us that the category of quasi-coherent sheaves \(\mathrm{Qch}(X)\) is equivalent to \(Mod(A)\). We can say this in the following way.

\begin{proposition}
     Let \(\mathcal{C}\) be a Grothendieck category with a projective generator \(\mathcal{O}\) and assume that \(\mathcal{O}\) is small, that is, \(\hom_{\mathcal{C}}(\mathcal{O},-)\) commutes with direct sums. Then \(\mathcal{C}\sim \Mod(A)\) for some ring \(A=\mathrm{End}(\mathcal{O})\).
\end{proposition}

Now consider the projective case. Let \(X\) be a projective scheme and \(\mathcal{L}\) be an ample line bundle over \(X\). For any \(n\geq 0\), define \(B_n=\Gamma(X,\mathcal{L}^n)\) with multiplication coming from the following composition 
\[\Gamma(X,\mathcal{L}^m)\otimes \Gamma(X,\mathcal{L}^n)\rightarrow \Gamma(X,\mathcal{L}^m\otimes \mathcal{L}^n)\rightarrow \Gamma(X,\mathcal{L}^{m+n}).\]
Then \(B=\oplus_{n\geq 0}B_n\) is a graded ring. Similarly, if \(\mathcal{M}\in \mathrm{Qch}(X)\), then we obtain a graded module 
\[\Gamma_h(\mathcal{M})=\oplus_{n\geq 0}\Gamma(X,\mathcal{M}\otimes \mathcal{L}^n).\]
\(\Gamma_h(-)\) defines a functor from \(\mathrm{Qch}(X)\) to \(\mathrm{Gr}(B)\) and we can project it to \(\mathrm{QGr}(B)\). This yields a functor 
\[\bar{\Gamma}_h:\mathrm{Qch}(X)\rightarrow \mathrm{QGr}(B).\]

\begin{theorem}[Serre's theorem]
  \begin{enumerate}[(i)]
    \item Let \(\mathcal{L}\) be an ample line bundle on a projective scheme \(X\). Then the functor \(\bar{\Gamma}_h\) defines an equivalence of categories between \(\mathrm{Qch}(X)\) and \(\mathrm{QGr}(B)\).
    \item Conversely, if \(R\) is a commutative connected graded \(k\)-algebra that is generated by \(R_1\), then there exists a line bundle \(\mathcal{L}\) over \(X=\proj R\) such that \(R=B(X,\mathcal{L})\) up to a finite dimensional vector space. Once again, \(\mathrm{QGr}(R)\simeq \mathrm{Qch}(X)\).
  \end{enumerate}
\end{theorem}

\begin{remark}
     It is essential that \(R\) is generated in degree 1. 
\end{remark}